\clearpage
\fancyhead[L]{\small VESTA. CERES. COMET.}
\SourceLeaf{313}\begin{center}
VESTA. CERES. COMET.
\end{center}

\par\vspace{4pt}\hrule\vspace{4pt}
\begin{center}
\fontsize{10}{12}\selectfont Göttingische gelehrte Anzeigen. Stück 180. Seite 1793..1795. 1808 November 10.
\end{center}
\par\vspace{4pt}\hrule\vspace{4pt}

Von Hrn. Prof. \textsc{Bode} sind uns folgende schätzbare, auf der Berliner Sternwarte angestellte, Beobachtungen der \emph{Vesta} mitgetheilt worden:


\begin{center}\fontsize{10}{12}\selectfont\setlength{\tabcolsep}{5pt}\renewcommand{\arraystretch}{1.15}\begin{tabular}{lrrr}
1808 & \shortstack{Mittlere Zeit\\in Berlin} & \shortstack{Scheinb. gerade\\Aufsteigung} & \shortstack{Scheinb. südl.\\Abweichung}\\
\hline
Aug. 24 & $11^{\text{u}}\;16^{\text{m}}\;16^{\text{s}}$ & $354^\circ\;34'\;59''$ & $13^\circ\;50'\;22''$\\
Sept. 1 & $12\;48\;0.7$ & $352\;57\;27$ & $14\;49\;41$\\
5 & $11\;8\;0$ & $352\;5\;20$ & $15\;20\;11$\\
13 & $11\;50\;1.0$ & $350\;15\;10$ & $16\;15\;59$\\
15 & $11\;40\;21.7$ & $349\;48\;2$ & $16\;28\;2$\\
20 & $11\;16\;26.7$ & $348\;42\;5$ & $16\;54\;35$\\
Oct. 6 & $10\;1\;55.2$ & $345\;49\;31$ & $17\;42\;7$\\
9 & $9\;48\;36.0$ & $345\;26\;47$ & $17\;44\;4$\\
14 & $9\;26\;57.3$ & $344\;56\;58$ & $17\;42\;30$\\
\end{tabular}\end{center}


\par\indent Alle diese Beobachtungen, die erste und dritte ausgenommen, sind am Mauerquadranten gemacht.

\par\indent Die \emph{Ceres} blieb am Mauerquadranten stets unsichtbar: folgende drei Beobachtungen sind mit dem
Kreismicrometer gemacht:


\begin{center}\fontsize{10}{12}\selectfont\setlength{\tabcolsep}{5pt}\renewcommand{\arraystretch}{1.15}\begin{tabular}{lrrr}
1808 & \shortstack{Mittlere Zeit\\in Berlin} & \shortstack{Scheinb. gerade\\Aufsteigung} & \shortstack{Scheinb. südl.\\Abweichung}\\
\hline
Aug. 29 & $10^{\text{u}}\;52^{\text{m}}\;49^{\text{s}}$ & $314^\circ\;45'\;57''$ & $30^\circ\;34'\;28''$\\
Sept. 16 & $7\;46\;36$ & $312\;26\;31$ & $30\;39\;58$\\
17 & $9\;12\;43$ & $312\;22\;24$ & $30\;40\;4$\\
\end{tabular}\end{center}


\par\indent Die folgenden, zu Petersburg von den Herren \textsc{Schubert} und \textsc{Wisniewsky} angestellten, Beobachtungen des letzten Cometen, welche uns gleichfalls durch die Güte des Hrn. Prof. \textsc{Bode} mitgetheilt sind
erhalten dadurch ein besonderes Interesse, dass sie die spätesten sind, die irgendwo haben gemacht werden können. Bei den letzten war der Comet so schwach, dass nur ein sehr scharfes Auge ihn im Fernrohre erkennen konnte. Die Beobachtungen sind mit dem Kreismicrometer gemacht, aber bei der Reduction hat man Refraction, Nutation und Aberration nicht in Betracht gezogen. Die verglichenen
Sterne hat man nicht bezeichnet.


\begin{center}\fontsize{10}{12}\selectfont\setlength{\tabcolsep}{5pt}\renewcommand{\arraystretch}{1.15}\begin{tabular}{lrrr}
1808 & \shortstack{Mittlere Zeit\\in Petersburg} & \shortstack{Gerade\\Aufsteigung} & \shortstack{Nordliche\\Abweichung}\\
\hline
Jan. 15 & $6^{\text{u}}\;49^{\text{m}}\;30^{\text{s}}$ & $335^\circ\;50'\;2''$ & $47^\circ\;47'\;27''$\\
26 & $7\;2\;25$ & $345\;33\;24$ & $48\;4\;6$\\
Febr. 15 & $8\;45\;44$ & $2\;36\;58$ & $48\;44\;20$\\
25 & $9\;17\;3$ & $9\;21\;23$ & $48\;50\;48$\\
März 1 & $9\;53\;5$ & $12\;33\;34$ & $48\;51\;59$\\
22 & $9\;24\;29$ & $23\;6\;39$ & $48\;48\;25$\\
27 & $10\;54\;20$ & $25\;48\;51$ & $48\;53\;49$\\
\end{tabular}\end{center}


\par\indent Wir fügen bei dieser Gelegenheit auch noch die auf der hiesigen Sternwarte in diesem Jahre angestellten Beobachtungen der \emph{Vesta} bei; die erste, welche überhaupt die früheste in diesem Jahre ge\SourceLeaf{314}machte zu sein scheint, ist bereits im 107. St. dieser Blätter [S. 303. d. Bandes] angezeigt. Vom 10. August an sind die Beobachtungen am Mauerquadranten gemacht; zufällige Hindernisse haben die Fortsetzung im October unterbrochen.


\begin{center}\fontsize{10}{12}\selectfont\setlength{\tabcolsep}{5pt}\renewcommand{\arraystretch}{1.15}\begin{tabular}{lrrr}
1808 & \shortstack{Mittlere Zeit\\in Göttingen} & \shortstack{Scheinbare gerade\\Aufsteigung} & \shortstack{Scheinbare südl.\\Abweichung}\\
\hline
Juni 22 & $13^{\text{u}}\;46^{\text{m}}\;25^{\text{s}}$ & $353^\circ\;12'\;17''$ & $9^\circ\;12'\;16''$\\
Juli  1 & $14\;38\;57$ & $354\;28\;3$ & $9\;9\;39$\\
2 & $12\;50\;2$ & $355\;7\;7$ & \\
13 & $13\;22\;37$ & $356\;35\;27$ & $9\;27\;40$\\
Aug. 10 & $14\;28\;53.7$ & $356\;36\;16.5$ & $11\;52\;50.3$\\
12 & $14\;20\;22.3$ & $356\;23\;22.3$ & $12\;8\;5.3$\\
14 & $14\;11\;23.4$ & $356\;9\;3.7$ & $12\;23\;25.7$\\
15 & $14\;6\;54.8$ & $356\;1\;6.0$ & $12\;31\;21.2$\\
22 & $13\;35\;4.7$ & $354\;56\;6.3$ & $13\;28\;20.7$\\
23 & $13\;30\;28.0$ & $354\;45\;24.1$ & $13\;36\;10.4$\\
29 & $13\;2\;11.5$ & $353\;35\;29.2$ & $14\;25\;25.8$\\
30 & $12\;57\;26.4$ & $353\;23\;12.9$ & $14\;33\;32.6$\\
Sept. 5 & $12\;28\;37.0$ & $352\;4\;33.3$ & $15\;20\;32.5$\\
8 & $12\;14\;6.6$ & $351\;23\;46.5$ & $15\;42\;22.3$\\
11 & $11\;59\;34.6$ & $350\;43\;7.3$ & $16\;3\;7.1$\\
12 & $11\;54\;42.9$ & $350\;28\;40.5$ & $16\;9\;34.9$\\
21 & $11\;11\;22.7$ & $348\;29\;10.5$ & $16\;59\;51.1$\\
\end{tabular}\end{center}


\par\indent Einige dieser Beobachtungen sind von Hrn. \textsc{Tiarks}, welcher sich bei uns den mathematischen Wissenschaften mit glücklichem Erfolge widmet, und im astronomischen Calcul bereits viele Fertigkeit besitzt, mit denjenigen Elementen verglichen, welche wir im 107. St. dieser Anz. [S. 303 d. B.] mitgetheilt haben, woraus folgende Differenzen hervorgegangen sind:


\begin{center}\fontsize{10}{12}\selectfont\setlength{\tabcolsep}{5pt}\renewcommand{\arraystretch}{1.15}\begin{tabular}{lrr}
 & \multicolumn{2}{c}{Differenz}\\
 & in der Länge & in der Breite\\\hline
Aug. 23 & $+4'\;7''4$ & $+43''0$\\
29 & $+4\;5.6$ & $+50.6$\\
30 & $+3\;53.9$ & $+45.0$\\
Sept. 5 & $+4\;4.5$ & $+45.2$\\
8 & $+4\;11.3$ & $+54.0$\\
11 & $+3\;31.6$ & $+38.3$\\
12 & $+4\;7.7$ & $+59.6$\\
21 & $+3\;37.0$ & $+55.1$\\
\end{tabular}\end{center}


\Needspace{16\baselineskip}
\fancyhead[L]{\small VESTA.}
\begin{center}
VESTA.

\par\vspace{4pt}\hrule\vspace{4pt}

\fontsize{9}{11}\selectfont Göttingische gelehrte Anzeigen. Stück 55. Seite 537..540. 1809. April 8.

\par\vspace{4pt}\hrule\vspace{4pt}
\end{center}

Wir haben in einigen frühern Blättern unsrer Gel. Anz. die ersten, vom Juni bis September des vorigen Jahres auf der hiesigen Sternwarte, wie auch die in Berlin von Hrn. Prof. \textsc{Bode} angestellten, Beobachtungen der Vesta mitgetheilt (1808 S. 1793) [Nov. 10. S. 308 d. B.]: vor kurzem hat Hr. Prof. \textsc{Gauss} auch die auf der kaiserl. Sternwarte zu Paris gemachten Beobachtungen dieses Planeten von Hrn.
\SourceLeaf{315}\textsc{Bouvard} mitgetheilt erhalten, durch deren frühere Bekanntmachung wir den Astronomen einen Dienst zu erweisen glauben.

\begin{center}
Beobachtungen der Vesta in Paris:
\end{center}


\begin{center}\fontsize{10}{12}\selectfont\setlength{\tabcolsep}{5pt}\renewcommand{\arraystretch}{1.15}\begin{tabular}{lrrr}
1808 & \shortstack{Mittlere Zeit\\in Paris} & \shortstack{Scheinbare gerade\\Aufsteigung} & \shortstack{Scheinbare südl.\\Abweichung}\\
\hline
Aug. 26 & $13^{\text{u}}\;16^{\text{m}}\;23^{\text{s}}5$ & $354^\circ\;11'\;28''80$ & $14^\circ\;1'\;20''3$\\
28 & $13\;6\;56.1$ & $353\;47\;30.30$ & $14\;17\;43.4$\\
Sept. 16 & $11\;35\;23.2$ & $349\;34\;9.45$ & $16\;34\;8.7$\\
19 & $11\;20\;57.3$ & $348\;54\;30.00$ & $16\;50\;16.0$\\
22 & $11\;6\;37.0$ & $348\;16\;15.00$ & $17\;4\;24.0$\\
24 & $10\;57\;7.7$ & $347\;51\;49.80$ & $17\;12\;48.5$\\
25 & $10\;52\;24.5$ & $347\;39\;55.35$ & $17\;16\;36.0$\\
26 & $10\;47\;42.1$ & $347\;28\;18.75$ & $17\;20\;10.0$\\
Oct.  4 & $10\;10\;48.6$ & $346\;6\;32.55$ & $17\;39\;50.5$\\
5 & $10\;6\;18.0$ & $345\;57\;50.25$ & $17\;41\;6.7$\\
6 & $10\;1\;48.7$ & $345\;49\;27.75$ & $17\;42\;20.0$\\
13 & $9\;31\;7.4$ & $345\;1\;51.75$ & $17\;43\;36.3$\\
\end{tabular}\end{center}


{\rightskip=0pt plus 1em \par\indent Aus diesen und den übrigen vorhandenen Beobachtungen hat Hr. Prof. \textsc{Gauss} den Gegenschein der Vesta bestimmt, den \emph{ersten}, der bisher beobachtet worden ist.\par}

\Needspace{7\baselineskip}
\begin{center}Opposition der Vesta:

1808 September 8. $7^{\text{u}}\;52^{\text{m}}\;57^{\text{s}}$ Mittl. Zeit in Göttingen.
\end{center}
\begin{center}\fontsize{10}{12}\selectfont\setlength{\tabcolsep}{5pt}\renewcommand{\arraystretch}{1.15}\begin{tabular}{lr}

Wahre Länge .\;.\;.\;.\;.\;.\;.\;.\;.\; & $345^\circ\;53'\;38''5$\\
Wahre geocentrische Breite .\;.\;.\;.\; & $11\;0\;25.8$ südl.\\
\end{tabular}\end{center}


\par\indent Auf der hiesigen Sternwarte wurden die Meridianbeobachtungen der Vesta bis zum November fortgesetzt, dann aber aufgegeben, da bei dem abnehmenden Lichte des Planeten die geringe Öffnung des Fernrohrs an unserm Mauerquadranten keine sehr genaue Bestimmung mehr erlaubte, und weniger genaue Beobachtungen aus dieser Zeit doch für die Verbesserung der Elemente von wenigem Nutzen gewesen sein würden. Erst gegen das Ende der diesmaligen Sichtbarkeit wurden desshalb die Beobachtungen mit dem Kreismicrometer wieder angefangen: das unbeständige Wetter hat zwar nur zwei Beobachtungen zu machen verstattet, die aber vorzüglich gut ausgefallen sind, und daher um so mehr hier
mitgetheilt zu werden verdienen, da sie wahrscheinlich die spätesten diesmal gemachten sind.


\begin{center}\fontsize{10}{12}\selectfont\setlength{\tabcolsep}{5pt}\renewcommand{\arraystretch}{1.15}\begin{tabular}{lrrr}
1809 & \shortstack{Mittlere Zeit\\in Göttingen} & \shortstack{Scheinbare gerade\\Aufsteig. der Vesta} & \shortstack{Scheinbare südl.\\Abweichung}\\
\hline
Febr. 8 & $7^{\text{u}}\;8^{\text{m}}\;25^{\text{s}}$ & $9^\circ\;26'\;51''3$ & $2^\circ\;9'\;28''8$\\
16 & $7\;16\;56$ & $12\;25\;20.0$ & $0\;41\;32.4$\\
\end{tabular}\end{center}


\par\indent Die im 107. Stück unsrer vorigjährigen Anzeigen [1808. Juli 4. S. 302 d. B.] abgedruckten Elemente weichen von diesen letzten Beobachtungen um 4 Min. in der Länge, und nur wenige Secunden in der Breite ab. Jene Elemente haben daher nur einer mässigen Verbesserung bedurft, um mit den Beobachtungen 1807, 1808 und 1809 in Übereinstimmung gebracht zu werden. Folgendes sind die neuen Elemente, welche Hr. Prof. \textsc{Gauss} neuerlich bestimmt hat.



\Needspace{20\baselineskip}
\SourceLeaf{316}\begin{center}
IV. \emph{Elemente der Vesta.}
\end{center}


\begin{center}\fontsize{10}{12}\selectfont\setlength{\tabcolsep}{5pt}\renewcommand{\arraystretch}{1.15}\begin{tabular}{lr}
\multicolumn{2}{c}{Epochen der mittlern Länge.}\\
\multicolumn{2}{c}{Meridian von Göttingen.}\\
1807 .\;.\;.\;.\;.\;.\;.\;.\;.\; & $168^\circ\;16'\;35''5$\\
1808 .\;.\;.\;.\;.\;.\;.\;.\;.\; & $267\;39\;28.6$\\
1809 .\;.\;.\;.\;.\;.\;.\;.\;.\; & $6\;46\;4.2$\\
1810 .\;.\;.\;.\;.\;.\;.\;.\;.\; & $105\;52\;39.7$\\
1811 .\;.\;.\;.\;.\;.\;.\;.\;\textsuperscript{.}\; & $204\;59\;15.3$\\
\end{tabular}\end{center}

\begin{center}\fontsize{9}{11}\selectfont\setlength{\tabcolsep}{5pt}\renewcommand{\arraystretch}{1.15}\begin{tabular}{lr}

Sonnennähe 1807, siderisch ruhend .\; & $249^\circ\;52'\;23''8$\\
Aufst. Knoten 1807, siderisch ruhend  & $103\;13\;11{,}2$\\
Neigung der Bahn .\;.\;.\;.\;.\;.\; & $7\;8\;18{,}8$\\
Tägliche mittl. tropische Bewegung .\;.\; & $977''5221$\\
Tropische Umlaufszeit .\;.\; & 1325 Tage 19 Stunden\\
Excentricität .\;.\;.\;.\;.\;.\;.\;.\;.\;.\; & $0{,}0887809$\\
Logarithm der halben grossen Axe .\;.\; & $0{,}3732940$\\
\end{tabular}\end{center}


\par\indent Die zweite Opposition der Vesta haben wir nach diesen Elementen den 1. Januar 1810 Nachmittags in $100^\circ\;43'$ Länge, und $0^\circ\;31'$ südlicher geocentrischer Breite zu erwarten.

\par\indent Vielleicht sind einigen unsrer Leser auch die Resultate über die verschiedene relative Lichtstärke dieses Planeten nicht uninteressant, welche Hr. Prof. \textsc{Gauss} bei diesen Rechnungen im Vorbeigehen mit entwickelt hat. Diejenige Lichtstärke als Einheit angenommen, welche der Planet in der Distanz 1 von der Erde und Sonne zeigen würde, war die Lichtstärke am Tage der ersten Entdeckung, den 29. März 1807, wo der Planet sich als ein Stern von 5..6. Grösse zeigte, $= 0,11815$; sie war damals schon im Abnehmen, und am 26. September 1807, wo die Mailänder Astronomen den Planeten zuletzt beobachteten, bis zu 0.02699 herabgesunken. Am 22. Juni 1808, wo der Planet auf unsrer Sternwarte zuerst als ein Stern 8. oder 7..8. Grösse wiedergefunden wurde, war die Lichtstärke schon wieder $= 0.05492$; der Planet wurde immer heller, und glich im September, wo seine Lichtstärke in der Opposition auf 0,09513 angewachsen war, einem Sterne 6. oder 6..7. Grösse. Jetzt nahm die Lichtstärke wieder ab, so dass sie am 16. Februar, wo doch der Planet noch reichlich die Helligkeit eines Sterns 9. Grösse hatte, auf 0,01629 gesunken war. In der nächsten Opposition wird sie nur bis auf 0,06435 anwachsen, und daher der Planet nur in dem Lichte eines Sterns 7. Grösse zu erwarten sein.
